%%%%%%%%%%%%%%%%%%%%%%%%%%%%%%%%%%%%%%%%%%%%%%%%%%%%%%%%%%%%%%%%%%%%%%
% Lớp 11 | Dãy số. Cấp số cộng. Cấp số nhân | Cấp số nhân
% Dạng 1D2?3-1: Câu hỏi lý thuyết
% Số câu: 8  (sinh tự động bởi tach_cau.py; từ nay thầy sửa trực tiếp file này)
%%%%%%%%%%%%%%%%%%%%%%%%%%%%%%%%%%%%%%%%%%%%%%%%%%%%%%%%%%%%%%%%%%%%%%

%%%%%%%%%%%%%%%% TRẮC NGHIỆM 4 PHƯƠNG ÁN %%%%%%%%%%%%%%%%

% nguồn: Ok_Toan11_TracNghiem_Lan2
\begin{ex}%[1D2N3-1]%[ma:d97cc288]
	Cho cấp số nhân $(u_n)$ có công bội $q$. Chọn hệ thức đúng trong  các hệ thức sau
	\choice 
		 {$u_k=\sqrt{u_{k+1}\cdot u_{k+2}}$}
		{$u_k=\dfrac{u_{k+1}+u_{k+2}}{2}$} 
		{\True $u_k=u_1\cdot  q^{k-1}$}
		{$u_k=u_1+(k-1)q$} 
	\loigiai{ 
	 Công thức số hạng tổng quát của cấp số nhân là $u_k=u_1\cdot q^{k-1}$.
	} 
\end{ex}

% nguồn: Ok_TracNghiemLop11_Lan1
\begin{ex}%[1D2N3-1]%[ma:90dd36cb]
	Cho cấp số nhân $(u_n)$ có số hạng đầu $u_1=-2$ và $u_6=486$. Công bội $q$  bằng
	\choice 
		 {$q=3$}
		{\True $q=-3$} 
		 {$q=\dfrac{1}{3}$}
		{$q=-\dfrac{1}{3}$} 
	 \loigiai{
		Ta có  $u_6=486\Leftrightarrow u_1 \cdot q^5=486 \Leftrightarrow (-2) \cdot q^5=486 \Rightarrow q^5=-243 \Leftrightarrow q=-3$.
	} 
\end{ex}

% nguồn: Ok_TracNghiemLop11_Lan1
\begin{ex}%[1D2N3-1]%[ma:856eae46]
	Cho $(u_n)$ là cấp số nhân lùi vô hạn có số hạng đầu và  công bội lần lượt là $u_1$ và $q$. Công thức nào sau đây dùng để tính tổng $S$ của cấp số nhân trên?
	 \choice
		 {$S=\dfrac{1-q}{u_1}$}
		 {$S=\dfrac{u_1}{1-q}$}
		{$S=\dfrac{q-1}{u_1}$} 
		 {\True $S=\dfrac{u_1}{1-q}$}
	\loigiai{ 
		Cho $(u_n)$ là cấp số nhân lùi vô hạn có số hạng đầu và công bội lần lượt là $u_1$ và  $q$. Công thức dùng để tính tổng $S$ là $S=\dfrac{u_1}{1-q}$.
	} 
\end{ex}

% nguồn: Ok_Toan11_TracNghiem_Lan2
\begin{ex}%[1D2H3-1]%[ma:13f3168a]
	Trong các  phát biểu sau, phát biểu nào là sai?
	\choice 
		{Dãy số có tất cả các số hạng bằng nhau là  một cấp số nhân}
		{Dãy số có tất cả  các số hạng bằng nhau là một cấp số cộng}
		{Một cấp số cộng  có công sai dương là một dãy số tăng}
		{\True Một cấp số cộng có công sai  dương là một dãy số dương}
	\loigiai{ 
		\begin{itemize}
			\item Dãy số có tất cả các số hạng bằng nhau là một cấp số nhân. Đúng vì dãy số đã cho là cấp số  nhân với công bội $q=1$.
			\item Dãy số có tất cả  các số hạng bằng nhau là một cấp số cộng. Đúng vì dãy số đã cho là cấp số cộng với công sai $d=0$.
			\item Một cấp số cộng có công sai dương là một dãy số tăng. Đúng vì dãy số đã  cho là cấp số cộng có công sai dương nên: $u_{n+1}-u_n=d > 0$ $\Rightarrow{u_{n+1}}>u_n$.
			\item Một cấp số cộng có công sai dương là một dãy số dương. Sai, ví dụ dãy $-5$; $-2$;  $1$; $3$; $\ldots$ là dãy số có $d=3 > 0$ nhưng không phải là dãy số dương.
	\end{itemize}}
\end{ex}

% nguồn: Ok_TracNghiemLop11_Lan1
\begin{ex}%[1D2H3-1]%[ma:6ebba398]
	Mệnh đề nào  dưới đây $\textbf{sai}$?
	\choice 
		{Dãy số có tất  cả số hạng bằng nhau là một cấp số nhân}
		{Dãy số có tất cả số hạng bằng nhau  là một cấp số cộng}
		{Một cấp số cộng có công sai dương là một  dãy số tăng}
		 {\True Một cấp số nhân có công bội $q>1$ là một dãy tăng}
	 \loigiai{
		\begin{enumerate}[$\bullet$]
			 \item Dãy số có tất cả số hạng bằng nhau là một cấp số nhân là mệnh đề đúng.\\
			Vì xét dãy số $\left(u_n\right)$ là một cấp số  nhân nên ta có $u_{n+1} = q \cdot u_n$ với khi $u_n \neq 0$, $q=1$ thì ta có $u_{n+1} = u_n$.
			\item Dãy số có tất cả số hạng bằng nhau là một  cấp số cộng là mệnh đề đúng.\\
			Vì xét dãy số $\left(u_n\right)$ là một cấp  số cộng nên ta có $u_{n+1} =u_n + d$ với khi $d=0$ thì ta có $u_{n+1} = u_n$.
			\item Một cấp số cộng có công sai dương là một dãy số tăng  là mệnh đề đúng.\\
			Ta xét dãy $\left(u_n\right)$ là một cấp số cộng có công sai $d>0$. Vì $u_{n+1} = u_n + d$ nên $u_{n+1}  - u_n = d>0$. Do đó dãy $\left(u_n\right)$ tăng.
			\item Một cấp số nhân có công bội $q>1$ là một dãy tăng là  mệnh đề$\textbf{sai}$.\\
			Ta xét  dãy $\left(u_n\right)$ là một cấp số cộng có công bội $q>1$. Vì $u_{n+1}=u_n\cdot q$ với $u_n \neq 0$, $q>1$. Khi đó $u_{n+1}-u_n=u_n \left(q-1\right)$.\\
			Khi $u_n<0$ thì $u_{n+1}-u_n <0$. Do đó $\left(u_n\right)$  là dãy số giảm.
		\end{enumerate}
	 }
\end{ex}

% nguồn: Ok_TracNghiemLop11_Lan1
\begin{ex}%[1D2H3-1]%[ma:a5a0b5c2]
	Cho cấp số nhân $(u_n)$ có số hạng đầu $u_1$ và công bội $q$. Số hạng tổng quát $u_n$ được xác định theo công  thức:
	\choice 
		{$u_n=u_1+(n-1)q$ với $n  \ge 2$}
		{$u_n=u_1+nq$ với $n \ge  2$}
		{$u_n=u_1 \cdot q^n$ với $n \ge  2$}
		{\True  $u_n=u_1 \cdot q^{n-1}$ với $n \ge 2$}
	 \loigiai{
		Cấp số nhân $(u_n)$ có  số hạng đầu $u_1$ và công bội $q$. Số hạng tổng quát $(u_n)$ được xác định theo công thức: $u_n=u_1 \cdot q^{n-1}$ với $n \ge 2$.
	} 
\end{ex}


%%%%%%%%%%%%%%%% ĐÚNG/SAI %%%%%%%%%%%%%%%%

% nguồn: Ok_dungsai_lop11_Lan2
\begin{ex}%[1D2H3-1]%[ma:edfda9b7]
	Cho cấp số nhân $\left(u_n\right)$ có số hạng đầu $u_1$ và công bội là $q$. Xét tính đúng sai của các  mệnh đề sau
	 \choiceTF
		{\True $u_n=u_1 \cdot q^{n-1}$,  $\forall n \geq 1$, $n \in \mathbb{N}$}
		{$u_n=\dfrac{u_{n-1} \cdot u_{n+1}}{2}$, $\forall n \geq 2$,  $n \in \mathbb{N}$}
		{\True $u_{n-1}\cdot u_{n+1}=u_n^2$, $\forall n \geq 2$, $n  \in \mathbb{N}$}
		{$u_1+u_2+\ldots+u_n=\dfrac{n\left(u_1+u_n\right)}{2}$, $n \in \mathbb{N}^*$} 
	 \loigiai{
		\begin{itemchoice}
			 \itemch \textbf{Đúng} \\
			Ta có $u_n=u_1 \cdot q^{n-1}$, $\forall n \geq 1$, $n \in  \mathbb{N}$.
			\itemch  \textbf{Sai} \\
			Ta có $u_{n-1}\cdot u_{n+1}=u_n^2$, $\forall n \geq  2$, $n \in \mathbb{N}$
			\itemch \textbf{Đúng}  \\
			Ta có $u_{n-1}\cdot u_{n+1}=u_n^2$, $\forall n \geq 2$, $n  \in \mathbb{N}$.
			\itemch \textbf{Sai}  \\
			Ta có $u_1+u_2+\ldots+u_n=\dfrac{u_1\left(1-q^n\right)}{1-q}$, $q \neq  1$.
		\end{itemchoice}
	} 
\end{ex}


%%%%%%%%%%%%%%%% TRẢ LỜI NGẮN %%%%%%%%%%%%%%%%

% nguồn: Ok_TraLoiNgan_Lop11_Lan1
\begin{ex}%[1D2T3-1]%[ma:e7a1e6c9]
	Bạn Lan thả quả bóng cao su từ độ cao $12$ mét theo phương thẳng đứng. Mỗi khi chạm đất nó lại nảy lên  theo phương thẳng đứng với độ cao bằng $\dfrac{2}{3}$ độ cao trước đó. Tính tổng quãng đường bóng đi được đến khi bóng dừng hẳn.
	 \shortans{60}
	\loigiai{ 
		Các quãng đường khi bóng đi xuống tạo thành một cấp số  nhân có $u_1=12$ và $q=\dfrac{2}{3}$.\\
		Tổng các quãng đường khi bóng  đi xuống là $S=\dfrac{u_1}{1-q}=\dfrac{12}{\left(1-\dfrac{2}{3}\right)}=36$ mét.\\
		Vậy tổng quãng đường bóng đi được (cả lên và xuống) đến khi bóng dừng hẳn  là
		$$2 S-12=2\cdot 36-12=60  \text{ (mét).}$$
	} 
\end{ex}
